\section{Evaluating the first application}
\label{main:evidence}

The framework makes three predictions about the first application, each naming a link: changing the source removes the transitions between layers; the failures that remain belong to the geometry map; and decomposition repairs them up to each piece's hull gap. A computational demonstration tests them and also measures the distribution control that separating order from shape is meant to provide. Its scope is stated in Section~\ref{sec:introduction}.

The implementation computes the semi-discrete map of Section~\ref{main:semidiscrete} with the solver of Geogram \cite{levy2015numerical,kitagawa2019convergence}, on parts meshed into tetrahedra \cite{si2015tetgen}, solving coarse to fine. The five parts are a cylinder, a sphere, a solid torus, an L-bracket with a reentrant edge, and a complex part consisting of a block with several holes and a boss (genus three). Every run, at orders $m=4,5,6$ ($4096$ to $262\,144$ waypoints), met the equal-mass condition to a relative tolerance of $10^{-6}$. A stop is counted wherever a path would have to be executed as travel: for the layered raster baseline, each travel move; for the transported path, each maximal \emph{exterior run} of segments that leave the part. Steps longer than five times the mean are reported separately as \emph{jumps}. Appendix~\ref{sec:implementation} gives the meshing, solver settings, and acceptance checks; Appendix~\ref{sec:experiments-setup} defines the metrics and the layered baseline; and Table~\ref{tab:experiments-geometries} in Appendix~\ref{sec:experiments-geometries} reports all fifteen runs.

\paragraph{A connected source removes the stops on convex parts.}
Figure~\ref{fig:layer-stops} shows the mechanism on a cube at order two. The decisive count is the number of stops between layers, which Proposition~\ref{prop:topology-transitions} forces. On the cylinder, a layered raster makes 13, 27, and 57 of them at orders 4, 5, and 6, roughly doubling with the number of layers, while the transported path has no layers to stop between and no exterior run at any tested order (Table~\ref{tab:main-stops}). That zero follows by construction on a convex part; beyond it, the path has no jump, and its longest step shrinks with the order. Changing the source thus removes a limit that grows with every refinement of the layered chain.

\begin{figure}[tb]
  \centering
  \includegraphicsorbox{figures/cube_layer_interruptions.pdf}{0.85\textwidth}{5.5cm}
  \caption{Layer changes on a cube at order two (64 waypoints). (a) Layered raster on the same cell centres: dashed red segments mark the three layer changes, each requiring a stop. (b) Computed transported path: it changes height throughout, so there is no layer boundary to stop at, and it has no exterior run. Colour indicates traversal time, from light (start, circle) to dark (end, square). Stop counts under refinement are in Table~\ref{tab:main-stops}.}
  \label{fig:layer-stops}
\end{figure}

% Written by python -m mapchain tables from results/evaluation/raster and results/evaluation/whole_part.
\begin{papertable}
  \begin{papertabular}{lrrr}
    \toprule
    & \multicolumn{2}{c}{Layered raster} & Transported path \\
    \cmidrule(lr){2-3}\cmidrule(lr){4-4}
    Part & Between layers & All travel & Exterior runs \\
    \midrule
    Cylinder & 13 / 27 / 57 & 97 / 419 / 2029 & 0 / 0 / 0 \\
    Sphere & 13 / 27 / 57 & 101 / 371 / 1533 & 0 / 0 / 0 \\
    Torus & 5 / 11 / 23 & 69 / 271 / 1163 & 1 / 53 / 241 \\
    L-bracket & 3 / 9 / 19 & 3 / 19 / 19 & 53 / 224 / 836 \\
    Complex part & 5 / 11 / 23 & 37 / 155 / 586 & 95 / 323 / 1103 \\
    \bottomrule
  \end{papertabular}
  \caption{Stops at Hilbert orders 4 / 5 / 6. Most raster travel lies within layers, where a better in-layer planner could remove some of it; only the stops between layers are forced.}
  \label{tab:main-stops}
\end{papertable}


\paragraph{The distribution is controlled cell by cell.}
Every accepted path assigns each waypoint the same target mass, so each waypoint stands for one Hilbert time block. On the convex parts the computable bound $d^{(m)}$ roughly halves with each order, from $0.120$ to $0.065$ and $0.035$ on the cylinder, and the covering radius follows. This is consistent with the $2^{-\beta m}$ term of \eqref{eq:main-bounds}, though three orders do not determine $\beta$.

\paragraph{On non-convex parts, the failures lie in the geometry map and persist under refinement.}
As the framework predicts, the source stays connected and the failures come from non-convex clipped cells (Figure~\ref{fig:containment-repair}). On the L-bracket, cells near the reentrant edge straddle the empty region between the arms and their centroids lie in it. On the torus at order four, every waypoint is inside the part, but one segment joins two interior centroids across the hole, which is also how most exterior segments of the complex part arise; Appendix~\ref{sec:experiments-resolution} separates the two mechanisms. Refinement lowers the fraction of exterior segments on the L-bracket and the complex part, but not the deepest excursions (Table~\ref{tab:experiments-geometries}): cells clipped by holes and reentrant edges stay large, and on the L-bracket $d^{(m)}$ stays near $0.85$, because refinement does not make the target convex. The cost is real: the torus has fewer exterior runs than the raster has travel moves, but the L-bracket and complex part have more (Table~\ref{tab:main-stops}). This persistence motivates decomposition.

\begin{figure}[tb]
  \centering
  \includegraphicsorbox{figures/nonconvex_containment_decomposition.pdf}{\textwidth}{7cm}
  \caption{Containment failures and their repair, computed at order four (plan views; top row L-bracket, bottom row torus; red marks the path outside the part). Left: whole-part paths (4096 waypoints). Middle: close-up of the boxed region; thick segments join consecutive waypoints (dots). On the L-bracket, an exterior run passes through a waypoint in the empty region beside the reentrant corner. On the torus, the single exterior segment joins two interior waypoints across the hole. Right: piece paths, with one colour per piece: the two convex boxes of the L-bracket at order four, and eight automatic torus pieces at order three (4096 waypoints in total, the whole-part budget; Appendix~\ref{sec:experiments-torus-figure}). No piece-path segment leaves the part.}
  \label{fig:containment-repair}
\end{figure}

\paragraph{Decomposition restores containment up to each piece's hull gap.}
We apply the same traversal and solver independently to the two boxes of the L-bracket and to an automatic partition of the torus into eight pieces; all ten local solves pass. The results match the two kinds of excursion of Section~\ref{main:decomposed}. Against each \emph{assigned piece}, the bound holds: the two convex boxes give full containment, as Proposition~\ref{prop:sd-hull-gap} guarantees (Figure~\ref{fig:containment-repair}, top right), and the automatic torus pieces, which are not exactly convex, have segments that leave their own piece, but every excursion stays below that piece's computed hull-gap bound (Table~\ref{tab:decomposition}). Against the \emph{whole object}, these crossings enter neighbouring pieces, not empty space: at the same budget of $32\,768$ waypoints, the whole-part path has 99 segments leaving the torus, while the piece paths have none, with similar coverage (Table~\ref{tab:torus-whole-vs-pieces}). These are sampled measurements with numerical hull-gap bounds, not containment certificates; Appendix~\ref{sec:experiments-pieces-independent} gives the partition procedure, the hull-gap evaluation, and the same-budget protocol.

% Written by python -m mapchain tables from results/evaluation/pieces/torus_whole_vs_pieces.json.
\begin{papertable}
  \begin{papertabular}{lrrrrrrr}
    \toprule
    Path / domain & Inside & Ext.\ seg. & Ext.\ length & Depth & Coverage & CV & Radius \\
    \midrule
    Whole / torus & 99.98\% & 0.30\% & 0.51\% & $0.0317$ & $0.981$ & $0.392$ & $0.0244$ \\
    Pieces / torus & 100.00\% & 0.00\% & 0.00\% & $0.0000$ & $0.978$ & $0.405$ & $0.0200$ \\
    Pieces / assigned & 98.40\% & 5.21\% & 3.11\% & $0.0232$ & -- & -- & -- \\
    \bottomrule
  \end{papertabular}
  \caption{Torus at the same budget: one whole-part path and eight piece paths, $32\,768$ waypoints each. \emph{Torus} measures against the whole object, \emph{assigned} against each path's own piece. ``Inside'' is the percentage of waypoints inside the evaluation domain; ``Ext.\ seg.'' and ``Ext.\ length'' are the percentages of exterior segments and exterior path length. Depth is the deepest sampled excursion; radius is the sampled covering radius. Distances are in normalised units; coverage and the coefficient of variation (CV) are defined in Appendix~\ref{sec:experiments-setup}. See Appendix~\ref{sec:experiments-torus-comparison} for the comparison protocol.}
  \label{tab:torus-whole-vs-pieces}
\end{papertable}


Decomposition thus trades a growing count for a fixed one: it brings back $N$ source components, independent of resolution, in exchange for containment up to each piece's hull gap. The first application kept the usual order of construction, geometry first and machine second. The second changes that order, because the next source of forced interruptions lies in the machine.
\FloatBarrier
